Automatic speech recognition (ASR) systems achieve strong performance in offline settings \citep{radford2023whisper,liu2025voxtral}, where the entire audio input is available before transcription begins. However, many real-world applications—such as voice assistants, live captioning, and interactive speech interfaces—require transcriptions to be produced in real time while audio is streaming, under strict latency constraints. Bridging the gap between offline transcription quality and real-time streaming remains a central challenge in speech recognition \citep{graves2012rnnt,zeghidour2025streaming}.

A common approach to streaming adapts offline models by processing audio in short chunks as it arrives \citep{machacek2023streaming}. While effective at moderate latencies, this strategy has a fundamental limitation: offline models are typically trained with access to bidirectional acoustic context (and often full-sequence conditioning), whereas a streaming system must emit tokens before future audio is available. This training--inference mismatch becomes increasingly severe as latency is reduced, and often leads to degraded accuracy in low-delay regimes and out-of-distribution settings.

\looseness=-1 Native streaming architectures address this by reformulating the learning problem so that each output token is predicted using only past inputs and a bounded amount of lookahead, making latency a tunable constraint. This requires (i) an explicit alignment between input audio and output text (e.g., word- or frame-level alignments), and (ii) an architecture that processes new audio incrementally as it arrives. A common instantiation is a neural transducer (RNN-T) with a streaming encoder that limits right context through chunking, memory, and caching \citep{graves2012rnnt,Shi2020EmformerEM, chen2021developing, noroozi2024nemotron}. Delayed Streams Modeling (DSM) \citep{zeghidour2025streaming} follows the same alignment-based principle, but replaces the transducer with a decoder-only model over aligned audio and text streams, enabling simpler designs that leverage pre-trained language decoders. DSM approaches offline accuracy at high delay settings. However, achieving offline-level performance at sub-second latency—particularly in multilingual and multi-domain settings—has remained an open challenge.

We address this challenge by introducing \emph{Voxtral Realtime}, a 4B parameter natively streaming ASR model that supports 13 languages. % Our architecture builds on DSM, extending it with a continuous causal encoder and adaptive delay conditioning so a single model can operate across streaming delays. 
Concretely, our contributions are:
\begin{itemize}
    \item A causal audio encoder trained from scratch with modern architectural choices (RMSNorm, SwiGLU, RoPE, sliding window attention).%, avoiding the mismatch introduced by adapting bidirectional offline encoders.
    \item An adaptive RMS-Norm (Ada RMS-Norm) mechanism in the decoder, enabling a single model to operate at any delay that is a multiple of 80\,ms.
    \item Pretraining at scale on a large-scale dataset spanning 13 languages, enabling robust generalization across languages and domains.
\end{itemize}

\looseness=-1 At a delay of 480\,ms, Voxtral Realtime achieves performance competitive with Whisper \citep{radford2023whisper} and ElevenLabs Scribe v2 Realtime \citep{scribev22025}. At higher delay settings (e.g., 960\,ms), it matches or surpasses strong offline baselines such as Voxtral Mini Transcribe V2 on several English and multilingual benchmarks \citep{voxtral22026}. These results demonstrate that offline-level transcription quality can be achieved within a fully streaming framework at sub-second latency.

We release the resulting model as open weights under the Apache 2.0 license. The remainder of this report details the model architecture, training and inference methodology, and empirical evaluations that support these findings.
